\chapter{Weg 1: Ubuntu auf einem eigenen Rechner installieren}

Dieses Kapitel ist für Sie, wenn Sie \textbf{keine USB-SSD} nutzen, sondern einen Rechner dauerhaft für NOMAD einrichten (zum Beispiel ein altes Notebook). Alles, was auf dem Rechner ist, wird \textbf{gelöscht}. Haben Sie die USB-SSD aus Kapitel~2 benutzt, überspringen Sie dieses und das nächste Kapitel und machen in Kapitel~5 weiter.

\hinweis{Die Bilder zeigen die Installation von Ubuntu 26.04 LTS in einer Testumgebung. Auf Ihrem Rechner sehen Sie dieselben Fenster, nur mit Ihren Namen und Plattengrößen.}

\section{Ubuntu-Stick erstellen}
Laden Sie die Datei „Ubuntu Desktop 26.04 LTS“ von \texttt{ubuntu.com/download/desktop} herunter (etwa 6~GB) und schreiben Sie sie mit balenaEtcher auf einen USB-Stick ab 8~GB, genau wie in Kapitel~2 beschrieben. Starten Sie den Rechner vom Stick (Startmenü, Tasten siehe Kapitel~2).

\section{Installation}
\bild[0.75]{02-sprache-de}{Sprache wählen: \menue{Deutsch}.}
\bild[0.75]{04-tastatur}{Tastaturbelegung: \menue{Deutsch}.}
\bild[0.75]{05-netzwerk}{Internetverbindung herstellen (Kabel oder WLAN). Ohne Internet geht die Installation nicht weiter.}
\bild[0.75]{07-art-der-installation}{Art der Installation: die normale Installation wählen.}
\bild[0.75]{10-festplatte}{Festplatte: „Festplatte löschen und Ubuntu installieren“. \textbf{Das löscht alles auf dieser Platte.}}
\bild[0.75]{12-konto-ausgefuellt}{Benutzerkonto: Name, Rechnername und ein Passwort, das Sie sich merken. Tragen Sie es in Kapitel~9 ein.}
\bild[0.75]{14-zusammenfassung}{Zusammenfassung: Prüfen und auf \menue{Installieren} klicken.}
\bild[0.75]{16-installation-fertig}{Wenn die Installation fertig ist, auf \menue{Jetzt neu starten} klicken. Den Stick danach abziehen, wenn der Rechner dazu auffordert.}
\bild[0.75]{17-anmeldung}{Mit dem Passwort aus dem Benutzerkonto anmelden.}

Nach der Anmeldung folgen einige Willkommens-Fenster. Sie können alle mit \menue{Weiter} oder \menue{Überspringen} beantworten. Zusätzliche Programme (Ubuntu~Pro) brauchen Sie nicht.

\merke{Prüfen Sie, dass der Rechner im Internet ist, bevor Sie weitermachen: Öffnen Sie den Browser (Firefox) und rufen Sie eine beliebige Seite auf.}
